% !TEX program = xelatex
% ============================================================
%  Arabic Independent Contractor Agreement Template — Nibras Center
%  عقد مقاول مستقل / متعاون حر
%  Compile with: xelatex main.tex (twice)
% ============================================================

\documentclass[12pt,a4paper]{article}

\usepackage[a4paper,margin=2.5cm,top=2.5cm,bottom=2.5cm,headheight=15pt]{geometry}
\usepackage{array}
\usepackage{booktabs}
\usepackage{tabularx}
\usepackage{enumitem}
\usepackage{float}
\usepackage{titlesec}
\usepackage{fancyhdr}
\usepackage{setspace}

\usepackage{bidi}
\usepackage[no-math]{fontspec}
\usepackage[quiet]{polyglossia}

\setmainfont[Scale=1]{NotoKufiArabic-Regular.ttf}
\newfontfamily\arabicfont[Script=Arabic,Scale=0.95]{NotoKufiArabic-Regular.ttf}
\newfontfamily\englishfont[Scale=0.95]{TeX Gyre Termes}

\setdefaultlanguage[calendar=gregorian,hijricorrection=1]{arabic}
\setotherlanguage[variant=british]{english}

\usepackage[breaklinks,hidelinks]{hyperref}
\hypersetup{pdftitle={عقد مقاول مستقل},pdfauthor={الشركة}}

\titleformat{\section}{\large\bfseries}{\thesection}{0.7em}{}
\titlespacing*{\section}{0pt}{1.2ex plus 0.3ex}{0.8ex plus 0.2ex}

\pagestyle{fancy}
\fancyhf{}
\fancyhead[R]{\small\itshape عقد مقاول مستقل}
\fancyhead[L]{\small اسم الشركة}
\fancyfoot[C]{\small\thepage}
\renewcommand{\headrulewidth}{0.3pt}

\newcolumntype{L}[1]{>{\raggedright\arraybackslash}p{#1}}
\newcolumntype{C}[1]{>{\centering\arraybackslash}p{#1}}

\setstretch{1.5}

\begin{document}

\begin{center}
{\LARGE\bfseries عقد مقاول مستقل}\\[0.5em]
{\large (\LR{Independent Contractor Agreement})}\\[1em]
\end{center}

\vspace{1em}

% ============================================================
% PARTIES
% ============================================================
\noindent\textbf{الأطراف:}

\begin{enumerate}
    \item \textbf{العميل (الشركة):} \rule{8cm}{0.4pt}، رقم السجل التجاري: \rule{4cm}{0.4pt}، المقر: \rule{6cm}{0.4pt}.

    \item \textbf{المقاول (المتعاون الحر):} \rule{8cm}{0.4pt}، رقم الهوية/الجواز: \rule{4cm}{0.4pt}، العنوان: \rule{6cm}{0.4pt}.
\end{enumerate}

\vspace{1em}

% ============================================================
% 1. SCOPE OF WORK
% ============================================================
\section{نطاق العمل}
\label{sec:scope}
\begin{enumerate}
    \item \textbf{وصف العمل:} \rule{10cm}{0.4pt}
    \item \textbf{المخرجات المطلوبة:}
    \begin{itemize}
        \item \rule{10cm}{0.4pt}
        \item \rule{10cm}{0.4pt}
        \item \rule{10cm}{0.4pt}
    \end{itemize}
    \item \textbf{المعايير:} يجب أن تكون المخرجات وفقاً للمعايير المتعارف عليها في الصناعة.
    \item يجوز تعديل نطاق العمل باتفاق مكتوب بين الطرفين.
\end{enumerate}

% ============================================================
% 2. TIMELINE
% ============================================================
\section{الجدول الزمني}
\label{sec:timeline}

\begin{table}[H]
\centering
\renewcommand{\arraystretch}{1.5}
\begin{tabularx}{\textwidth}{C{1cm} L{6cm} C{3cm} X}
\toprule
\textbf{رقم} & \textbf{المخرج} & \textbf{تاريخ التسليم} & \textbf{ملاحظات} \\
\midrule
1 & --- & --- & --- \\
2 & --- & --- & --- \\
3 & --- & --- & --- \\
\bottomrule
\end{tabularx}
\end{table}

% ============================================================
% 3. PAYMENT
% ============================================================
\section{الأتعاف والمدفوعات}
\label{sec:payment}

\begin{table}[H]
\centering
\renewcommand{\arraystretch}{1.5}
\begin{tabularx}{\textwidth}{L{6cm} C{3cm} X}
\toprule
\textbf{البند} & \textbf{المبلغ (\LR{USD})} & \textbf{موعد الدفع} \\
\midrule
دفعة أولى (عند التوقيع) & --- & عند التوقيع \\
دفعة ثانية (منتصف المشروع) & --- & عند تسليم المخرج 2 \\
دفعة نهائية (عند التسليم) & --- & عند التسليم النهائي \\
\midrule
\textbf{الإجمالي} & \textbf{---} & \\
\bottomrule
\end{tabularx}
\end{table}

\begin{itemize}
    \item تُدفع الفواتير خلال \rule{2cm}{0.4pt} يوماً من استلامها.
    \item المسؤول عن الضرائب: المقاول مسؤول عن ضرائبه الخاصة.
    \item لا تُدفع مصاريف إضافية إلا بموافقة مسبقة من العميل.
\end{itemize}

% ============================================================
% 4. CONTRACTOR STATUS
% ============================================================
\section{صفة المقاول}
\label{sec:status}
\begin{enumerate}
    \item \textbf{المقاول ليس موظفاً} لدى العميل، ولا تنطبق عليه أحكام عقد العمل.
    \item لا يحق للمقاول الحصول على مزايا الموظفين (تأمين صحي، إجازات مدفوعة، تقاعد).
    \item يتحمل المقاول مسؤولية ضرائبه وتأمينه الخاص.
    \item يحق للمقاول تحديد ساعات ومكان عمله، شريطة الالتزام بالجدول الزمني.
    \item يحق للمقاول العمل لدى عملاء آخرين، شريطة عدم تعارض المصالح.
\end{enumerate}

% ============================================================
% 5. IP
% ============================================================
\section{الملكية الفكرية}
\label{sec:ip}
\begin{enumerate}
    \item تنتقل جميع حقوق الملكية الفكرية في المخرجات إلى العميل فور تسليمها ودفع الأتعاف.
    \item يوقّع المقاول اتفاقية تنازل منفصلة عن الملكية الفكرية.
    \item لا يحق للمقاول استخدام المخرجات لأغراض شخصية أو إعادة بيعها.
    \item الأدوات والقوالب التي يستخدمها المقاول قبل العقد تبقى ملكاً له، ويمنح العميل ترخيصاً لاستخدامها.
\end{enumerate}

% ============================================================
% 6. CONFIDENTIALITY
% ============================================================
\section{السرية}
\label{sec:confidentiality}
يلتزم المقاول بالحفاظ على سرية جميع معلومات العميل، وفقاً لاتفاقية عدم الإفصاح (\LR{NDA}) الموقّعة بين الطرفين. يستمر هذا الالتزام لمدة \rule{2cm}{0.4pt} سنوات بعد انتهاء العقد.

% ============================================================
% 7. WARRANTIES
% ============================================================
\section{الضمانات}
\label{sec:warranties}
\begin{enumerate}
    \item يضمن المقاول أن المخرجات أصلية ولا تنتهك حقوق ملكية فكرية لأطراف ثالثة.
    \item يضمن أن العمل سيكون خالياً من العيوب الجوهرية لمدة \rule{2cm}{0.4pt} يوماً بعد التسليم.
    \item يصلح المقاول أي عيوب تكتشف خلال فترة الضمان دون مقابل.
\end{enumerate}

% ============================================================
% 8. TERMINATION
% ============================================================
\section{إنهاء العقد}
\label{sec:termination}
\begin{enumerate}
    \item يجوز لأي طرف إنهاء العقد بإشعار مسبق \rule{2cm}{0.4pt} يوماً.
    \item يجوز للعميل الإنهاء فوراً «لسبب» (إخلال جسيم، تأخر متكرر).
    \item عند الإنهاء، يُدفع للمقاول مقابل العمل المنجز حتى تاريخ الإنهاء.
    \item يسلّم المقاول جميع المواد والوثائق المتعلقة بالمشروع.
\end{enumerate}

% ============================================================
% 9. LIABILITY
% ============================================================
\section{الحدّ من المسؤولية}
\label{sec:liability}
\begin{enumerate}
    \item لا تتجاوز مسؤولية المقاول إجمالي الأتعاف المدفوعة.
    \item لا يكون المقاول مسؤولاً عن الأضرار غير المباشرة.
    \item يلتزم المقاول بالحصول على تأمين المسؤولية المهنية.
\end{enumerate}

% ============================================================
% 10. GOVERNING LAW
% ============================================================
\section{القانون الحاكم}
\label{sec:law}
تُفسَّر هذه الاتفاقية وفقاً لقوانين \rule{6cm}{0.4pt}. تُحلّ النزاعات عن طريق التحكيم في \rule{4cm}{0.4pt}.

% ============================================================
% SIGNATURES
% ============================================================
\vspace{2em}
\section*{التوقيعات}

\begin{table}[H]
\centering
\renewcommand{\arraystretch}{2.5}
\begin{tabularx}{\textwidth}{L{5cm} X X}
\toprule
& \textbf{العميل} & \textbf{المقاول} \\
\midrule
\textbf{الاسم} & --- & --- \\
\textbf{الصفة} & --- & --- \\
\textbf{التوقيع} & \rule{4cm}{0.4pt} & \rule{4cm}{0.4pt} \\
\textbf{التاريخ} & \rule{4cm}{0.4pt} & \rule{4cm}{0.4pt} \\
\bottomrule
\end{tabularx}
\end{table}

\end{document}
